% =========================================================
% NANO-experience — Results
% =========================================================
%
% TODO – FIB patterning result (before DRIE)
%   - Describe the resist pattern obtained after development (ask Mario for images)
%   - Were the holes well-defined? Any proximity effects or column milling artefacts?
%
% TODO – Size comparison: NANO vs. MICRO
%   - MICRO features: µm-scale (limited by UV wavelength ~365–405 nm and diffraction)
%   - NANO features: ~50–150 nm (limited by Ga+ spot size, no diffraction issue)
%   - Key point: FIB uses ions not photons → no diffraction limit in the classical sense
%   - Discuss: ion straggle / scattering in resist vs. optical proximity effects in lithography
%
% TODO – SEM results (from DRIE experience, connected here)
%   - After ICP-DRIE: SEM images of the Si nanopillars
%   - Measure and report: pillar diameter, height, pitch
%   - Verticality: measure sidewall angle from vertical (ideally 90°)
%   - Angulation: any tilt of the pillar axis? Cause: beam tilt or non-uniform etch
%   - Compare observed geometry with target (150 nm diameter nanopillar)
%
% TODO – Discussion points required by the slides
%   1. Detailed description of FIB process + parameters (refer to M&M section)
%      Special focus on resist reticulation mechanism (cross-linking by Ga+ secondary e-)
%   2. Why smaller features than MICRO: wavelength/beam size argument
%   3. Verticality and angulation values from SEM — are they within spec?
%      Comment on what causes deviation (etch selectivity, mask erosion, DRIE loading)
%
% NOTE: Full structural characterisation (SEM images, angle measurements) comes from
% the DRIE experience. Cross-reference that section here.
